\documentclass[10pt,a4paper]{amsart}
\usepackage[margin=19mm]{geometry}
\usepackage{amsmath,amssymb,mathtools}
\usepackage[scr=rsfs,cal=euler]{mathalfa}
\usepackage{xfrac}
\usepackage[unicode,hidelinks,bookmarks=false]{hyperref}
\newcommand{\ie}{i.e.\ }
\newcommand{\cf}{cf.\ }
\newcommand{\resp}{respectively\ }
\newcommand{\Subsection}{\subsection}
\providecommand{\nobreakdash}{\nobreak\hbox{-}}
\setlength{\emergencystretch}{2em}
\multlinegap0pt
\usepackage{bm}
\usepackage{bbm}
\usepackage{dsfont}
\usepackage{xspace}
\usepackage{cancel}
\usepackage{mathtools}
\usepackage{amsthm}
\makeatletter
\@ifundefined{lemma}{\newtheorem{lemma}{Lemma}}{}
\@ifundefined{theorem}{\newtheorem{theorem}{Theorem}}{}
\makeatother
\newcommand{\E}{\mathbb{E}}
\newcommand{\Tr}{{\rm Tr}}
\title[Neural Thermodynamics: Entropic Forces in Deep and Universal]{Neural Thermodynamics: Entropic Forces in Deep and Universal Representation Learning}
\author{Liu Ziyin, Yizhou Xu, Isaac Chuang}
\date{Source: arXiv:2505.12387v4 (CC BY 4.0)}
\begin{document}
\maketitle

\noindent\emph{Source and licence.} Neural Thermodynamics: Entropic Forces in Deep and Universal Representation Learning. Version arXiv:2505.12387v4 (CC BY 4.0), distributed under the Creative Commons Attribution 4.0 International licence, as recorded in the arXiv licence metadata for this exact version and shown on the abstract page https://arxiv.org/abs/2505.12387v4. Full rights notice: licenses/arxiv-2505.12387-CC-BY-4.0.txt.\par\smallskip

\noindent\textit{Editorial scope.} Counted unit: the Theorem whose source label is \texttt{theo:deep\_linear} in the article (source0.tex lines 1092--1106, source label \texttt{theo:deep\_linear}), delivered together with the article's own complete original proof of it (source0.tex lines 1108--1179, one \texttt{proof} environment of 72 lines). No mathematics is written, repaired or re-derived: statement and proof are the article's own text line for line. Required context retained uncounted: eq: free energy (lines 185--188); theo:attention (lines 418--425); lemma:decomposition (lines 1236--1247). Every definition, display and label of those lines is retained unchanged. Omitted from this package and not redistributed: 1--184, 189--417, 426--1091, 1107--1107, 1180--1235, 1248--1564. Nothing inside a retained excerpt is omitted or altered: every retained body is delivered exactly as the article's lines are written, so no omission receipt of any kind accompanies this package. Citations: the retained excerpts carry no citation command at all, so no bibliography entry of the article is transcribed and no reference is redistributed. Reference adaptation: none was needed and none was made; every reference inside the delivered unit resolves inside it, so no unresolved reference or citation is produced. Printed slips kept literal: no printed word, symbol, hypothesis or number of the counted statement or of its proof was altered. Layout and class: the article's own class is \texttt{article} with packages including neurips\_2025, amsmath, amsthm, verbatim, amssymb, amsfonts, amscd, graphicx, enumitem, graphics, centernot, authblk, hyperref, footmisc; the wrapper is the lane's pinned \texttt{amsart} editorial wrapper at 10pt on A4, which restates the theorem-like environments the retained text needs and adds the article's own macro definitions that the retained text needs (\texttt{\textbackslash E}, \texttt{\textbackslash Tr}), with extra wrapper packages (bm, bbm, dsfont, xspace, cancel, mathtools). The delivered numbering is the unit's own. Assets: no retained span contains a figure environment, a graphics-inclusion command, a PDF-page inclusion, a table or a TikZ picture; no image file is copied into this package, the package declares no asset dependency and no image file is redistributed with it. Licence: version arXiv:2505.12387v4 of this article is distributed under the Creative Commons Attribution 4.0 International licence, as recorded in the arXiv licence metadata for this exact version and shown on the abstract page https://arxiv.org/abs/2505.12387v4. A full rights notice is delivered at licenses/arxiv-2505.12387-CC-BY-4.0.txt. Characters: every retained excerpt is pure ASCII, so the delivered engine, XeLaTeX with its default Latin Modern text font, reports no missing character.
\begin{align}
\label{eq: free energy}
F_{\eta,\gamma}(\theta)&=\underbrace{\mathbb{E}_x\ell(x,\theta)}_{\text{learning, symmetry}}+\underbrace{\gamma||\theta||^2}_{\text{regularization}} +\underbrace{\frac{1}{4}\mathbb{E}_\mathcal{B}||\sqrt{\Lambda}\mathbb{E}_{x\in\mathcal{B}}  \nabla \ell(x,\theta)||^2 }_{\text{effective entropy due to discretization error, noise, $:=S(\theta)$}} +  O(||\Lambda||^2).
\end{align}

\begin{theorem}
\label{theo:attention}
(Gradient Alignment) For all local minimum of Eq.\eqref{eq: free energy}, we have
\begin{equation}
\eta\E[G_W^T G_W-G_U G_U^T]=4\gamma(W^TW-UU^T).
\end{equation}

\end{theorem}

\begin{lemma}
\label{lemma:decomposition}
Suppose that matrices $W_1,W_2,\cdots,W_D$ satisfy
\begin{equation}
W_{i+1}^TW_{i+2}^T...W_D^TW_D...W_{i+2}W_{i+1}=\lambda_iW_iW_{i-1}...W_1W_1^T...W_{i-1}^TW_i
\end{equation}
for some $\lambda_i>0$ and $i=1,2,\cdots,D-1$, then we have write the SVD of $W_1,W_2,\cdots,W_D$ as
\begin{equation}
W_D=U_D\Sigma_DU_{D-1}^T,\ W_{D-1}=U_{D-1}\Sigma_{D-1}U_{D-2}^T,\cdots,\ W_1=U_1\Sigma_1U_0,
\end{equation}
where $\Sigma_D,\cdots,\Sigma_1$ are the singular values.
\end{lemma}

\begin{theorem}
\label{theo:deep_linear}
Let $V'=\sqrt{\Sigma_\epsilon}V\sqrt{\Sigma_x}$ such that $V'=\tilde{U}S'\tilde{V}$ is its SVD and $\text{rank}(V')=d$. Assume that every layer has more than $d$ hidden units. Then if $\gamma=0$ and $\eta=0^+$, at any global minimum of \eqref{eq: free energy}, we have
\begin{equation}
\sqrt{\Sigma_\epsilon}M_1W_{D}=\tilde{U}\Sigma_DU_{D-1}^T,\ W_i=U_i\Sigma_iU_{i-1}^T,\ W_1M_2M_3\sqrt{\Sigma_x}=U_1\Sigma_1\tilde{V},
\end{equation}
for $i=2,\cdots,D-1$, where $U_i$ are arbitrary matrices satisfying $U_i^TU_i=I_{d\times d}$, and $\Sigma_x=\mathbb{E}[xx^T]$, $\Sigma_\epsilon=\mathbb{E}[\epsilon\epsilon^T]$. Moreover,
\begin{equation}
\begin{aligned}
&\Sigma_1=(\Tr S')^{-\frac{D-2}{2D}}\frac{\Tr[M_2M_3\Sigma_xM_3^TM_2^T]^{\frac{D-1}{2D}}}{\Tr[M_1^T\Sigma_\epsilon M_1]^{\frac{1}{2D}}}\sqrt{S'},\\
&\Sigma_D=(\Tr S')^{-\frac{D-2}{2D}}\frac{\Tr[M_1^T\Sigma_\epsilon M_1]^{\frac{D-1}{2D}}}{\Tr[M_2M_3\Sigma_xM_3^TM_2^T]^{\frac{1}{2D}}}\sqrt{S'},\\
&\Sigma_i=(\Tr S')^{1/D}(\Tr[M_1^T\Sigma_\epsilon M_1]\Tr[M_2M_3\Sigma_xM_3^TM_2^T])^{-\frac{1}{2D}}I_d.
\end{aligned}
\end{equation}
\end{theorem}

\begin{proof}
Consider two consecutive layers $W_i$ and $W_{i+1}$. Using Theorem \ref{theo:attention}, we have
\begin{equation}
\eta\E[G_{W_{i+1}}^T G_{W_{i+1}}-G_{W_i} G_{W_i}^T]=0.
\end{equation}
By the MSE loss $\ell(x,y)=||y-M_1W_D\cdots W_1M_2M_3x||^2$, this gives
\begin{equation}
W_ih_i\mathbb{E}[||\xi_{i+1}^T\tilde{r}||^2\tilde{x}\tilde{x}^T]h_i^TW_i^T=W_{i+1}^T\xi_{i+1}^T\mathbb{E}[||h_i\tilde{x}||^2\tilde{r}\tilde{r}^T]\xi_{i+1}W_{i+1},
\end{equation}
where $\tilde{x}=\E_{x\in\mathcal{B}}x$ and $\tilde{r}:=\E_{x\in\mathcal{B}}[y-M_1W_D\cdots W_1M_2M_3x]$ satisfy $\E\tilde{x}=\E\tilde{r}=0$, and thus $\E\tilde{x}\tilde{x}^T=\frac{\Sigma_x}{|\mathcal{B}|}$ and $\E\tilde{\epsilon}\tilde{\epsilon}^T=\frac{\Sigma_\epsilon}{|\mathcal{B}|}$. We use the fact that at the global minimum we have $M_1W_D\cdots W_1M_2M_3=V$, and thus $y-M_1W_D\cdots W_1M_2M_3x$ is independent of $x$. We denote $\xi_{i+1}:=M_1W_D\cdots W_{i+2}$, $h_i:=W_{i-1}\cdots W_1M_2M_3$ for $i=2,3,\cdots,D-2$, and $\xi_D:=M_1$, $h_1:=M_2M_3$.

Finally denote $W_1'=W_1M_2M_3\sqrt{\Sigma_x}$ and $W_D'=\sqrt{\Sigma_\epsilon}M_1W_D$, which gives $W_D'\cdots W_1'=V'$ and 
\begin{equation}
W_{i+1}^{\top} \frac{W_{i+2}^{\top} \cdots W_D^{\prime \top} W_D^{\prime} \cdots W_{i+2}}{\Tr \left[W_{i+2}^{\top} \cdots W_D^{\prime \top} W_D^{\prime} \cdots W_{i+2}\right]} W_{i+1}=W_i \frac{W_{i-1} \cdots W_1^{\prime} W_1^{\prime \top} \cdots W_{i-1}^{\top}}{\Tr \left[W_{i-1} \cdots W_1^{\prime} W_1^{\prime \top} \cdots W_{i-1}^{\top}\right]} W_i^{\top}
\label{eq:eqi}
\end{equation}
for $i=2,3,\cdots,D-2$. For $i=1$ we have
\begin{equation}
W_{2}^{\top} \frac{W_{3}^{\top} \cdots W_D^{\prime \top} W_D^{\prime} \cdots W_{3}}{\Tr \left[W_{3}^{\top} \cdots W_D^{\prime \top} W_D^{\prime} \cdots W_{3}\right]} W_{2}= \frac{W_1^{\prime} W_1^{\prime \top} }{\Tr \left[M_2M_3\Sigma_x M_2^\top M_3^\top\right]} 
\label{eq:eq1}
\end{equation}
and for $i=D-1$ we have
\begin{equation}
\frac{W_D^{\prime \top} W_D^{\prime}}{\Tr \left[M_1^T\Sigma_\epsilon M_1\right]}=W_{D-1} \frac{W_{D-2} \cdots W_1^{\prime} W_1^{\prime \top} \cdots W_{D-2}^{\top}}{\Tr \left[W_{D-2} \cdots W_1^{\prime} W_1^{\prime \top} \cdots W_{D-2}^{\top}\right]} W_{D-1}^{\top}.
\label{eq:eqD}
\end{equation}
Lemma \ref{lemma:decomposition} proves that we can decompose the matrices $W_1',W_2,\cdots,W_{D-1},W_D'$ as
\begin{equation}
W_D'=U_D\Sigma_DU_{D-1}^T,\ W_{D-1}=U_{D-1}\Sigma_{D-1}U_{D-2}^T,\cdots,\ W_1'=U_1\Sigma_1U_0.
\end{equation}

By minimizing \eqref{eq: free energy} at $\eta=0^+$, we need to minimize
\begin{equation}
\mathbb{E}||y-W_D\cdots W_1x||^2=\mathbb{E}||\epsilon||^2+||(V-W_D\cdots W_1)\sqrt{\Sigma_x}||^2.
\end{equation}
Thus we have
\begin{equation}
(V-W_D\cdots W_1)\sqrt{\Sigma_x}=0,
\end{equation}
which gives
\begin{equation}
W_D'\cdots W_1'=\sqrt{\Sigma_\epsilon}V\sqrt{\Sigma_x}=V'.
\end{equation}
Then we obtain $U_D=\tilde{U}$, $U_0=\tilde{V}$ and
\begin{equation}
\Sigma_D\Sigma_{D-1}\cdots\Sigma_1=S'.
\label{eq:S'}
\end{equation}
We can assume $\Sigma_D,\cdots,\Sigma_1\in\mathbb{R}^{d\times d}$ because their ranks are the same by \eqref{eq:eqi}, \eqref{eq:eq1} and \eqref{eq:eqD}.

\eqref{eq:eqi} gives
\begin{equation}
\frac{\Sigma_{i+1}^2\Sigma_{i+2}^2\cdots\Sigma_D^2}{\Tr [\Sigma_{i+2}^2\cdots\Sigma_D^2]}=\frac{\Sigma_1^2\cdots\Sigma_{i-1}^2\Sigma_i^2}{\Tr[\Sigma_1^2\cdots\Sigma_{i-1}^2]},
\end{equation}
and thus $\Sigma_i=cI_d$ for $i=2,3,\cdots,D-2$ and
\begin{equation}
\frac{\Sigma_1^2}{\Tr[\Sigma_1^2]}=\frac{\Sigma_D^2}{\Tr[\Sigma_D^2]}.
\label{eq:Sigma1D_1}
\end{equation}
\eqref{eq:eq1} and \eqref{eq:eqD} give
\begin{equation}
\frac{\Sigma_{2}^2\Sigma_{3}^2\cdots\Sigma_D^2}{\Tr [\Sigma_{3}^2\cdots\Sigma_D^2]}=\frac{\Sigma_1^2}{\Tr[M_2M_3\Sigma_xM_3^TM_2^T]},\ \frac{\Sigma_D^2}{\Tr[M_1^T\Sigma_\epsilon M_1]}=\frac{\Sigma_1^2\cdots\Sigma_{D-1}^2}{\Tr[\Sigma_1^2\cdots\Sigma_{D-2}^2]},
\end{equation}
and thus
\begin{equation}
c^2\frac{\Sigma_D^2}{\Tr[\Sigma_D^2]}=\frac{\Sigma_1^2}{\Tr[M_2M_3\Sigma_xM_3^TM_2^T]},\ \frac{\Sigma_D^2}{\Tr[M_1^T\Sigma_\epsilon M_1]}=c^2\frac{\Sigma_1^2}{\Tr[\Sigma_1^2]}.
\label{eq:Sigma1D_2}
\end{equation}
Combining \eqref{eq:S'}, \eqref{eq:Sigma1D_1} and \eqref{eq:Sigma1D_2}, we finish the proof.

% The following steps are the same as \cite[Theorem 5.4]{ziyin2024parameter}. We can verify that the solution satisfies $M_1[\Pi_{i=1}^{D_A}W_i^A]M_2M_3=V$, or $h_A^{D_A}(x)=Vx$ as expected. 
\end{proof}

\end{document}
